% !TEX root = ../main.tex
\section{Reglas del juego}\label{sec:reglas}

\subsection{Reglas de negocio implementadas}

Las reglas de negocio de la Tabla~\ref{tab:rn} se identifican con un código que se referencia desde los casos de uso y desde el modelo de dominio.

{\tablacompacta
\begin{longtable}{@{}L{1.5cm}L{14cm}@{}}
  \caption{Reglas de negocio del juego}\label{tab:rn}\\
  \toprule
  \cab{Código} & \cab{Regla} \\ \midrule \endfirsthead
  \toprule
  \cab{Código} & \cab{Regla} \\ \midrule \endhead
  \bottomrule \endlastfoot
  RN-01 & Los tipos de dominó disponibles son doble~6, con 28 fichas, y doble~9, con 55 fichas. \\ \filasep
  RN-02 & La partida tiene un jugador humano y de una a tres CPU, para un total de dos a cuatro jugadores. Siempre se juega de forma individual. \\ \filasep
  RN-03 & Cada jugador recibe 7 fichas en doble~6 y 10 fichas en doble~9. Las fichas restantes quedan disponibles para robar. \\ \filasep
  RN-04 & El límite de tiempo por turno puede ser sin límite, 10, 20 o 30 segundos, y aplica solo al jugador humano. Si el tiempo se agota, el jugador pierde el turno. \\ \filasep
  RN-05 & Abre la mano el jugador que tenga el doble más alto entre todos los jugadores, buscando en orden descendente desde el doble máximo del tipo de dominó hasta el doble~0. Si ningún jugador tiene dobles, abre quien tenga la ficha de mayor puntaje; en caso de empate, la que tenga el número individual más alto. \\ \filasep
  RN-06 & Una ficha solo puede colocarse en un extremo de la mesa con el que coincida uno de sus valores. \\ \filasep
  RN-07 & Si el jugador no tiene ficha para lanzar, debe robar hasta obtener una ficha jugable o hasta que no queden fichas para robar. Si no quedan fichas para robar y no tiene ficha jugable, pasa. \\ \filasep
  RN-08 & Gana la partida el primer jugador que se queda sin fichas. \\ \filasep
  RN-09 & El juego se cierra cuando no quedan fichas para robar y ningún jugador tiene ficha para lanzar. En ese caso se revelan las fichas de todos los jugadores para corroborar el resultado, y gana quien tenga la menor suma de puntos en sus fichas; en empate, quien tenga menos fichas; si el empate persiste, la partida termina empatada. Perder el turno por tiempo no cuenta para declarar el juego cerrado. \\ \filasep
  RN-10 & El jugador humano puede deshacer su última jugada. Al hacerlo se revierten también las jugadas de las CPU realizadas después de ella, hasta volver a su turno. \\ \filasep
  RN-11 & La partida en curso puede guardarse y cargarse posteriormente en el mismo estado. \\
\end{longtable}}

\subsection{Reparto y fichas para robar según la configuración}

La cantidad de fichas que quedan para robar depende de la combinación entre el tipo de dominó y el número de jugadores, como se resume en la Tabla~\ref{tab:robar}. Con doble~6 y cuatro jugadores se reparten las 28 fichas, de modo que no quedan fichas para robar y la regla RN-07 se reduce a pasar.

\begin{table}[H]
\caption{Fichas para robar por tipo de dominó y número de jugadores}\label{tab:robar}
\tablacompacta
\begin{tabularx}{0.75\textwidth}{@{}Y C{3.2cm}C{3.2cm}@{}}
  \toprule
  \cab{Jugadores} & \cab{Doble 6} & \cab{Doble 9} \\
  & {\footnotesize 7 fichas c/u} & {\footnotesize 10 fichas c/u} \\ \midrule
  Humano y 1 CPU & 14 & 35 \\
  Humano y 2 CPU & 7 & 25 \\
  Humano y 3 CPU & Ninguna & 15 \\
  \bottomrule
\end{tabularx}
\end{table}

\subsection{Reglas excluidas del alcance}

Las siguientes variantes quedan por fuera del proyecto. Las razones de cada exclusión se registran en la bitácora de decisiones (sección~\ref{sec:bitacora}).

\begin{enumerate}
  \item Juego en parejas.
  \item Multijugador en red.
  \item Doble~12 y otros tipos de dominó.
  \item Puntos acumulados entre manos y meta de puntos.
  \item Deshacer más de una jugada y rehacer.
  \item Bonificaciones especiales, como la capicúa.
\end{enumerate}
